% BHU PYQ -- Manually transcribed & error-corrected
% Paper: BCH-214 : Cost Accounting - I
% Exam: B.Com. (Hons.) Semester III Examination, 2013-14

\documentclass[11pt,a4paper]{article}
\usepackage[a4paper,top=2.5cm,bottom=2.5cm,left=2.8cm,right=2.8cm,headheight=14pt]{geometry}
\usepackage{amsmath,amssymb}
\usepackage[T1]{fontenc}
\usepackage[utf8]{inputenc}
\usepackage{lmodern,microtype,fancyhdr,enumitem,hyperref,lastpage,parskip}

\hypersetup{
    colorlinks=true,
    urlcolor=black,
    linkcolor=black,
    pdftitle={BCH-214: Cost Accounting - I -- BHU B.Com. (Hons.) Sem III 2013-14}
}

\pagestyle{fancy}
\fancyhf{}
\renewcommand{\headrulewidth}{0.4pt}
\renewcommand{\footrulewidth}{0.4pt}
\fancyhead[L]{\small\textit{Banaras Hindu University}}
\fancyhead[R]{\small\textit{BCH-214: Cost Accounting - I}}
\fancyfoot[C]{\small Page \thepage\ of \pageref{LastPage}}

\newlist{parts}{enumerate}{2}
\setlist[parts,1]{label=(\alph*),leftmargin=*,itemsep=5pt,topsep=3pt}
\setlist[parts,2]{label=(\roman*),leftmargin=*,itemsep=3pt,topsep=2pt}
\newcommand{\pts}[1]{\hfill\textit{[#1]}}

\begin{document}

\begin{center}
  {\large\bfseries BANARAS HINDU UNIVERSITY}\\[2pt]
  {\normalsize B.Com. (Hons.) --- Semester III Examination, 2013-14}\\[6pt]
  \rule{0.8\linewidth}{0.5pt}\\[6pt]
  {\large\bfseries Paper No. BCH-214: Cost Accounting - I}\\[4pt]
  {\normalsize Time: 3 Hours\hfill Maximum Marks: 70}
\end{center}

\rule{\linewidth}{0.4pt}

\smallskip
\noindent\textit{Instructions: Answer all questions. All questions carry equal marks (14 marks each).}

\bigskip

\noindent\textbf{Question 1.} \\
It is said, ``Cost Accounting is a system of foresight and not a postmortem examination, it turns losses into profits, speeds up activities and eliminates wastes''. Discuss in detail this Statement.\pts{14}

\centerline{\textbf{OR}}

What are the various methods of costing? Explain.\pts{14}

\medskip\rule{\linewidth}{0.2pt}

\noindent\textbf{Question 2.} \\
What is meant by overhead expenses? Discuss the different methods of allocating works on cost.\pts{14}

\centerline{\textbf{OR}}

Record the following transactions in the Stores Ledger Account of Udai Ltd.:\pts{14}
\begin{center}
\renewcommand{\arraystretch}{1.2}
\begin{tabular}{|l|l|r|r|}
\hline
\textbf{Date} & \textbf{Transaction} & \textbf{Units} & \textbf{Total Cost (Rs.)} \\ \hline
10.07.2012 & Receipt & 250 & 312.50 \\
21.08.2012 & Receipt & 100 & 130.00 \\
20.09.2012 & Receipt & 50 & 67.50 \\
05.10.2012 & Issue & 85 & -- \\
19.10.2012 & Receipt & 50 & 70.00 \\
25.10.2012 & Issue & 300 & -- \\
05.11.2012 & Issue & 40 & -- \\ \hline
\end{tabular}
\end{center}
The issues on 05.10.2012 and 25.10.2012 were priced on LIFO and FIFO basis respectively. As from 1.11.2012 it was decided to price the issues at weighted average price.

\medskip\rule{\linewidth}{0.2pt}

\noindent\textbf{Question 3.} \\
\begin{parts}
  \item What is material consumed? How is it calculated in unit costing?\pts{7}
  \item How would you do the valuation of cost of goods sold in unit costing?\pts{7}
\end{parts}

\centerline{\textbf{OR}}

From the following information, prepare a Reconciliation Statement and ascertain the profit/loss as per Financial Accounts:\pts{14}
\begin{itemize}[itemsep=2pt,topsep=2pt]
  \item Net loss as per costing records: Rs. 1,72,400
  \item Administrative overhead in excess in Cost A/c: Rs. 1,700
  \item Works overhead under recovered in Cost A/c: Rs. 3,120
  \item Depreciation in Financial Accounts: Rs. 11,400
  \item Depreciation in Cost Accounts: Rs. 12,700
  \item Interest received not included in Cost Accounts: Rs. 8,000
  \item Obsolescence loss charged in Financial Accounts: Rs. 5,700
  \item Income tax provided in Financial Accounts: Rs. 40,300
  \item Opening stock in Cost Accounts: Rs. 53,000
  \item Opening stock in Financial Accounts: Rs. 54,400
  \item Interest (Dr.) in Cost Accounts: Rs. 6,000
  \item Preliminary expenses written off: Rs. 800
  \item Closing stock in Cost Accounts: Rs. 52,000
  \item Closing stock in Financial Accounts: Rs. 49,600
  \item Provision for doubtful debts: Rs. 150
  \item Bank interest (Cr.) in Financial Accounts: Rs. 750
\end{itemize}

\medskip\rule{\linewidth}{0.2pt}

\noindent\textbf{Question 4.} \\
Differentiate between by-product, and joint products. How are by-products valued in Cost Accounting?\pts{14}

\centerline{\textbf{OR}}

In an oil refinery the product passes through three different processes. The following information is available for the month of January, 2013:\pts{14}
\begin{center}
\renewcommand{\arraystretch}{1.2}
\begin{tabular}{|l|r|r|r|}
\hline
\textbf{Particulars} & \textbf{Crushing (Rs.)} & \textbf{Refining (Rs.)} & \textbf{Finishing (Rs.)} \\ \hline
Raw materials (500 tons of Copra) & 2,25,000 & -- & -- \\
Wages & 8,000 & 5,900 & 5,875 \\
Power & 1,500 & 1,200 & 1,500 \\
Sundry materials & 300 & 1,700 & -- \\
Factory expenses & 600 & 1,000 & 950 \\ \hline
\end{tabular}
\end{center}
Cost of drums for storing finished oil was Rs. 21,025. 200 tons of cake were sold for Rs. 15,000 and 275 tons of crude oil were obtained. Sundry by-product of the crushing process fetched Rs. 900. By-product after refining the oil was sold for Rs. 900 (20 tons) and 250 tons of refined oil were obtained. 240 tons of finished oil were stored in drums and 10 tons were sold for Rs. 1,200. The establishment expenses for the period amounted to Rs. 3,500 which are to be charged to the three processes in proportion of 3:2:2.

Prepare the Process Accounts.

\medskip\rule{\linewidth}{0.2pt}

\noindent\textbf{Question 5.} \\
What are the essential features of contract costing? Explain the rules for the calculation of profit on incomplete contract.\pts{14}

\centerline{\textbf{OR}}

The authorised capital of Sudarshan Ltd. is Rs. 2,00,000 divided into 500 6\% preference shares of Rs. 100 each and 15,000 equity shares of Rs. 10 each. The company commenced contract work on January 1, 2012 and took a contract of Rs. 4,00,000 during the year. The Trial Balance of the company on 31st December, 2012 was as under:\pts{14}
\begin{center}
\renewcommand{\arraystretch}{1.2}
\begin{tabular}{|l|r|r|}
\hline
\textbf{Particulars} & \textbf{Dr. Balance (Rs.)} & \textbf{Cr. Balance (Rs.)} \\ \hline
500 6\% Preference shares of Rs. 100 each (Rs. 80 paid up) & -- & 40,000 \\
5,000 Equity shares of Rs. 10 each (Rs. 8 paid up) & -- & 40,000 \\
Sundry creditors & -- & 8,000 \\
Land and building & 34,000 & -- \\
Cash at bank & 9,000 & -- \\
Materials & 80,000 & -- \\
Plant & 15,000 & -- \\
Wages & 1,05,000 & -- \\
Expense & 5,000 & -- \\
Cash received (80\% of work certified) & -- & 1,60,000 \\ \hline
\textbf{Total} & \textbf{2,48,000} & \textbf{2,48,000} \\ \hline
\end{tabular}
\end{center}
Out of the plant and material sent on the contract site, plant worth Rs. 3,000 and material worth Rs. 2,400 were destroyed in an accident. Plant which originally cost Rs. 4,000 was returned to stores on 31st December, 2012. Material at site was Rs. 3,000 and the cost of uncertified work was Rs. 2,000. Depreciate plant @ 10\%.

Prepare Contract Account for the year 2012 and Balance Sheet as on 31st December, 2012.

\vfill
\rule{\linewidth}{0.4pt}
\begin{center}\small
\href{https://bhu-exam.vercel.app/}{Click here to visit our website}\\[4pt]
\href{https://me-aryan.vercel.app/}{Click here to visit our contributor}\\[4pt]
\href{https://quashan-me.vercel.app/}{Click here to visit our contributor}
\end{center}

\end{document}
